\subsection{Ошибка компиляции из-за отсутствующего заголовочного файла}

При добавлении функции \texttt{GetNetworkScanMutex()} в заголовочный файл \texttt{ui.h}
возникла ошибка компиляции в нескольких модулях. Функция возвращала
ссылку на \texttt{std::mutex}, но заголовочный файл \texttt{<mutex>} не был включён в
\texttt{ui.h}, что приводило к ошибке компиляции во всех модулях, которые включали
\texttt{ui.h}.

Информацию об ошибке можно рассмотреть ниже.

\begin{lstlisting}[style=CodeListing]
//psf/Home/Documents/university/5_term/CurseWork/file-manager/src/ui/layout.cpp: In function 'void UI::LayoutChildren(HWND, const UIControls&, const UIState&)':
//psf/Home/Documents/university/5_term/CurseWork/file-manager/src/ui/layout.cpp:23:15: error: 'mutex' in namespace 'std' does not name a type
   23 |         std::mutex& mtx = UI::GetNetworkScanMutex(hwnd);
      |               ^~~~~
In file included from //psf/Home/Documents/university/5_term/CurseWork/file-manager/src/ui/layout.cpp:10:
//psf/Home/Documents/university/5_term/CurseWork/file-manager/include/ui.h:50:8: note: 'std::mutex' has not been declared
   50 |     std::mutex& GetNetworkScanMutex(HWND hwnd);
      |        ^~~~~
\end{lstlisting}

Проблема была решена добавлением директивы \texttt{\#include <mutex>} в начало файла \texttt{ui.h}.
Это обеспечило объявление типа \texttt{std::mutex} во всех модулях, которые включают
заголовочный файл \texttt{ui.h}.

\begin{lstlisting}[style=CodeListing]
#include <mutex>

namespace UI {
    std::mutex& GetNetworkScanMutex(HWND hwnd);
}
\end{lstlisting}

После добавления \texttt{\#include <mutex>} все ошибки компиляции были устранены. Все модули
успешно компилируются, и функция \texttt{GetNetworkScanMutex()} корректно работает во всех местах
использования. Это подтвердило важность включения всех необходимых заголовочных файлов в
заголовочные файлы, которые используются другими модулями.
